\section*{Chapter 14 --- Random Variables and the Binomial Distribution}

\begin{solution}{exo:tle:randvar:1}
$\E(X) = -0.3 + 0 + 0.8 + 0.5 = 1$.
$\E(X^2) = 0.3 \times 1 + 0 + 0.4\times4 + 0.1\times25 = 4.4$, so by
König--Huygens $\V(X) = 4.4 - 1 = 3.4$ and
$\sigma(X) = \sqrt{3.4} \approx 1.84$.
\end{solution}

\begin{solution}{exo:tle:randvar:2}
\emph{1.} The $10$ questions are independent Bernoulli trials with
$p = \frac14$: $X \sim \mathcal B\left(10, \frac14\right)$,
$\E(X) = 2.5$,
$\sigma(X) = \sqrt{10 \times \frac14 \times \frac34} = \sqrt{1.875}
\approx 1.37$.

\emph{2.} $\P(X = 0) = \left(\frac34\right)^{10} \approx 0.056$;
\[
\P(X = 5) = \binom{10}{5}\left(\frac14\right)^5\left(\frac34\right)^5
\approx 0.058;
\qquad
\P(X \geq 1) = 1 - \left(\tfrac34\right)^{10} \approx 0.944 .
\]
\end{solution}

\begin{solution}{exo:tle:randvar:3}
Independent identical trials: $X \sim \mathcal B(400,\ 0.05)$, so
$\E(X) = 20$ claims and
$\sigma(X) = \sqrt{400 \times 0.05 \times 0.95} = \sqrt{19} \approx 4.4$.
\end{solution}

\begin{solution}{exo:tle:randvar:4}
The payment $P$ received satisfies
$\E(P) = \frac{5 + 6}{6} = \frac{11}{6}$, so the fair price is
$c = \frac{11}{6} \approx 1.83$ euros. The fair gain $G = P - \frac{11}6$
takes values $-\frac{11}{6}, \frac{19}{6}, \frac{25}{6}$ with probabilities
$\frac46, \frac16, \frac16$:
\[
\V(G) = \E(G^2)
= \frac{4 \times 121 + 361 + 625}{6 \times 36} = \frac{1470}{216}
= \frac{245}{36} \approx 6.8,
\qquad \sigma(G) \approx 2.6 .
\]
A fair game has zero \emph{average} gain but its outcomes still fluctuate:
fair is not riskless.
\end{solution}

\begin{solution}{exo:tle:randvar:5}
$X \sim \mathcal B(8,\ 0.7)$.
\[
\P(X = 6) = \binom86 (0.7)^6(0.3)^2 \approx 0.296 ;
\]
\[
\P(X \geq 6) = \P(6) + \P(7) + \P(8)
\approx 0.296 + 8(0.7)^7(0.3) + (0.7)^8
\approx 0.296 + 0.198 + 0.058 = 0.552 ;
\]
$\P(X \geq 1) = 1 - (0.3)^8 \approx 0.99993$. Mode:
$(n+1)p = 6.3$, so the most probable value is $k^* = 6$
(\cref{exo:tle:randvar:7}).
\end{solution}

\begin{solution}{exo:tle:randvar:6}
The number of passengers showing up is $X \sim \mathcal B(104,\ 0.9)$; the
flight is overbooked when $X \geq 101$:
\[
\P(X \geq 101) = \sum_{k=101}^{104}\binom{104}{k}(0.9)^k(0.1)^{104-k}
\approx 0.006 .
\]
Selling $4\%$ more tickets than seats causes an incident on only about
$0.6\%$ of flights --- the economics behind overbooking.
\end{solution}

\begin{solution}{exo:tle:randvar:7}
\[
\frac{\P(X = k+1)}{\P(X = k)}
= \frac{\binom{n}{k+1}}{\binom nk}\cdot\frac{p}{1-p}
= \frac{n - k}{k + 1}\cdot\frac{p}{1-p},
\]
using $\binom{n}{k+1} = \binom nk \frac{n-k}{k+1}$. This ratio is $\geq 1$
iff $(n-k)p \geq (k+1)(1-p)$ iff $np - k p \geq k - kp + 1 - p$ iff
$k \leq (n+1)p - 1$. So the probabilities increase strictly while
$k + 1 \leq (n+1)p$ and decrease afterwards: the maximum is attained at
$k^* = \floor{(n+1)p}$ (shared with $k^* - 1$ when $(n+1)p$ is an
integer).
\end{solution}

\begin{solution}{exo:tle:randvar:8}
\emph{1.} Heads first at toss $k$ means $k-1$ tails then heads:
probability $\left(\frac12\right)^{k-1}\cdot\frac12 = 2^{-k}$, for
$1 \leq k \leq 10$. Total winning probability
$\sum_{k=1}^{10} 2^{-k} = 1 - 2^{-10} = \frac{1023}{1024}$ (the game is
lost only on ten consecutive tails).

\emph{2.} Expected payoff:
\[
\sum_{k=1}^{10} 2^k \cdot 2^{-k} = \sum_{k=1}^{10} 1 = 10 \text{ euros}.
\]
Without the cap, the sum $\sum_{k\geq1} 1$ diverges: the expected payoff is
infinite, although the game almost always pays a small amount --- the
famous \emph{Saint Petersburg paradox}, showing that expectation alone does
not measure the value of a game.
\end{solution}
